\title{Byte Latent Transformer: Patches Scale Better Than Tokens}

\begin{abstract}
We present the Byte Latent Transformer ({\textsc BLT}), a novel byte-level large language model (LLM) framework that achieves, for the first time, parity with tokenization-based LLMs in terms of performance at scale, while offering substantial gains in inference speed and resilience. {\textsc BLT} processes input by encoding raw bytes into patches of variable lengths, which act as the model’s main computational elements. The segmentation of these patches is governed by the predicted entropy of subsequent bytes, thereby directing additional computational resources and representational power toward regions of higher data complexity. We conduct the first flop-constrained scaling analysis of byte-level models, training configurations with up to 8B parameters and ingesting 4T bytes of data. Our empirical findings indicate that it is indeed possible to scale models trained directly on bytes without relying on a predefined vocabulary. The model achieves more efficient training and inference by adaptively utilizing longer patches where input data exhibits higher predictability, alongside observed gains in reasoning capabilities and out-of-distribution generalization. At comparable inference budgets, {\textsc BLT} attains markedly better scaling performance than models reliant on tokenization, made possible by concurrently expanding both patch length and overall model capacity.
\end{abstract}

\section{Introduction}
\label{section:intro}

We present the Byte Latent Transformer~(\textbf{BLT}), an architecture that operates directly on raw byte sequences without relying on tokenization. BLT achieves parity with conventional tokenization-based models at scale, while also delivering marked gains in both efficiency and robustness (\S\ref{sec:robustness}).
Conventional large language models (llms) generally utilize an almost fully end-to-end training paradigm, with the notable exception of tokenization. This tokenization phase is a heuristic process in which bytes are organized into a predetermined set of discrete tokens. The resulting tokens shape the model’s compression strategy and introduce several biases, resulting in disadvantages such as sensitivity to domain or modality~\citep{dagangetting}, increased vulnerability to noisy inputs~(\S\ref{sec:robustness}), limited capacity for orthographic reasoning~\citep{edman2024cute}, and imbalanced treatment of languages~\citep{liang2023xlm,petrov2024language,limisiewicz2024myte}.

Historically, tokenization has been a critical component, as training llms directly on bytes incurs significant computational expense at scale because of increased sequence lengths~\citep{xue2022byt5}. Previous research has sought to address this challenge by introducing more computationally efficient forms of self-attention~\citep{el2020characterbert, clark2022canine} or by designing architectures that do not rely on attention mechanisms~\citep{wang2024mambabyte}~(\S\ref{section:related_work}). Yet, these strategies mainly benefit the training of \textit{small models}. When scaling up, the predominant computational bottleneck in Transformer models stems from the extensive feed-forward network layers applied to every byte, rather than the overhead associated with the attention mechanism.

We introduce a dynamic, trainable approach for aggregating bytes into \textit{patches}~(\S\ref{section:patching}), coupled with an innovative model architecture that integrates both byte-level and patch-level information. In contrast to tokenization, {\textsc BLT} dispenses with a predefined patch vocabulary. Instead, any sequence of bytes can be transformed into a latent patch embedding using compact, learned encoder and decoder components. Our findings demonstrate that this approach allows for \textit{greater} computational efficiency compared to traditional models reliant on tokenization.

Tokenization-based LLMs assign an equal computational budget to every token, prioritizing simplicity over optimal efficiency. However, because these tokens result from compression heuristics that do not necessarily align with the inherent complexity of each prediction, this strategy can be suboptimal. The core concept underpinning our architecture is that computational resources should be adaptively distributed according to demand. For instance, most word endings are straightforward, low-entropy predictions and thus do not require the computational power of a large transformer; such resources are better reserved for more complex predictions like generating the initial word of a new sentence. {\textsc BLT}'s design, as described in~(\S\ref{section:architecture}), operationalizes this principle by utilizing three transformer blocks: two compact byte-level \textit{local models} and one expansive global \textit{latent transformer}~(\autoref{fig:arch_high_level}). 

To achieve adaptive compute allocation, {\textsc BLT} employs a segmentation scheme based on the entropy associated with next-byte prediction. This results in the formation of context-sensitive groups of bytes that exhibit similar levels of information density.

We conduct the inaugural study on scaling byte-level models under controlled floating point operation (flop) budgets, evaluating architectures up to 8B parameters and 4T training bytes. Our results demonstrate that it is feasible to train large-scale, end-to-end models from raw bytes without relying on fixed-vocabulary tokenization schemes. Notably, {\textsc BLT} attains training flop-controlled performance on par with Llama 3, yet achieves up to a 50\% reduction in inference flops~(\S\ref{section:scaling})\footnote{We assess computational expenditure by tallying the total number of Floating Point OPerations (flops) executed.}. 

In addition, our findings indicate that operating directly on raw bytes leads to marked improvements in modeling infrequent or rare data patterns. {\textsc BLT} demonstrates increased robustness to input noise compared to traditional tokenizer-dependent models, and shows superior performance in tasks that demand detailed character-level comprehension, such as orthographic analysis, phonological evaluation, and low-resource machine translation~(\S\ref{sec:robustness}).

Furthermore, {\textsc BLT} models facilitate concurrent scaling of both model and patch sizes while adhering to a fixed inference flop budget. By employing larger patch sizes, the average compute per inference step decreases, enabling additional resources to be allocated towards expanding the global latent transformer, which is invoked less frequently as a result. Our flop-controlled scaling experiments during inference~(\autoref{fig:fixed_inference_scaling}) reveal consistently improved scaling behavior compared to models based on tokenization.

To summarize, this work offers several key contributions: 1) We propose {\textsc BLT}, a byte latent llm framework capable of adaptively distributing computational resources to enhance flop efficiency; 2) We provide evidence that, at scales up to 8B, our approach matches the training flop budget of Llama 3, with the flexibility to accept minor degradations in evaluation outcomes in exchange for up to 50\% improvements in flop efficiency; 3) {\textsc BLT} models introduce an alternative approach to scaling llms by enabling model size increases without exceeding a predetermined inference cost; 4) We show that {\textsc BLT} models possess greater resilience to input perturbations and demonstrate sensitivity to sub-word information typically overlooked by token-based llms. The codebase for both training and inference with {\textsc BLT} is available at~\url{https://github.com/facebookresearch/blt}.

\section{Patching: From Individual Bytes to Groups of Bytes}
\label{section:patching}

Dividing byte sequences into \textit{patches} enables {\textsc BLT} to flexibly assign computational resources according to contextual demands. As depicted in Figure~\ref{fig:patching_types}, there exist multiple strategies for partitioning bytes into patches. More precisely, a patching function $f_p$ maps an input sequence of bytes $\pmb{x}=\{x_i,|i=1,\ldots n\}$ of length $n$ into a shorter sequence of $m < n$ patches $\pmb{p}=\{p_j|j=1,\ldots,m\}$, with each $x_i$ being assigned to the set \{0,1\}; a value of 1 denotes the initiation of a new patch. For models that are either token-based or patch-based, the processing expense is dictated predominantly by the number of steps the central Transformer must execute. Specifically, in {\textsc BLT}, this equates to the number of patches produced by the patching function when encoding the data. Therefore, the typical length of a patch, referred to as the \textit{patch size}, serves as the principal determinant of the computational requirements for training and inference under a chosen patching function~(\S\ref{section:flops}).

In what follows, we describe three specific patching functions: fixed-size byte patching~(\S\ref{section:static-patch}), patching based on whitespace~(\S\ref{section:white-patch}), and patching dynamically informed by entropy estimations from a compact byte language model~(\S\ref{section:dyn-patch}). We conclude by exploring incremental patching and highlighting distinctions between patching and conventional tokenization approaches~(\S\ref{section:bpe-patch}).

\subsection{Strided Patching Every K Bytes}
\label{section:static-patch}

A simple and commonly used approach for organizing bytes is to segment them into uniform patches of size $k$, as in MegaByte~\citep{yu2023megabyte}. Employing a fixed stride facilitates both model training and inference, offers an intuitive means to adjust average patch size, and thus allows for straightforward control over computational complexity. Nevertheless, this patching approach presents notable drawbacks. Primarily, computational resources are not adaptively distributed to regions where they are most required: transformer step $j$ might be wasted, for instance, if it is solely predicting whitespace in programming code, or conversely, inadequate computation might be allotted to bytes that are highly information-dense, such as those representing mathematical expressions. Additionally, this strategy can produce inconsistent and non-contextual segmentation of comparable byte sequences, resulting in scenarios where identical words may be divided differently.

\subsection{Space Patching}
\label{section:white-patch}

\citet{slagle2024spacebyte} introduces a straightforward yet impactful enhancement to the strided patching approach by initiating new patches at every occurrence of a space-like byte\footnote{
    Space-like bytes are defined as any byte that is not a latin character, digit, or utf-8 continuation byte. In addition, each patch must contain at least one non space-like byte. }, which frequently serve as natural demarcations of linguistic units across multiple languages. Within the Space patching framework, an additional latent transformer step (i.e., increased computational effort) is dedicated to processing each word. This approach ensures that the patching strategy remains consistent for words across different sequences and directs computational resources toward more challenging prediction points, which often occur after spaces. As an illustration, generating the initial byte in the response to the prompt ``Who composed the Magic Flute?\uline{\hspace{.6em}}'' presents a greater challenge than producing subsequent bytes after ``M''; this is because the initial character considerably narrows the pool of plausible continuations, rendering the completion ``Mozart'' relatively straightforward to predict. Nevertheless, space patching exhibits limitations in its applicability across all languages and domains, and, crucially, it lacks the ability to flexibly adjust patch sizes. In the subsequent section, we propose a novel patching strategy that leverages the observation that initial bytes of words are generally more challenging to predict, while also introducing a mechanism for dynamically managing patch size.

\subsection{Entropy Patching: Using Next-Byte Entropies from a Small Byte LM}
\label{section:dyn-patch}

In place of traditional rule-based heuristics like whitespace detection, our methodology leverages a data-driven framework to pinpoint next-byte predictions with significant uncertainty. We propose \textit{entropy patching}, a technique that utilizes entropy measurements to determine the boundaries of each patch.

To implement this, we first train a compact byte-level auto-regressive language model using the {\textsc BLT} training corpus, then calculate the entropy for the subsequent byte under the LM-induced distribution $p_{e}$ over the byte vocabulary $\mathcal{V}$:

\begin{equation}
H(x_i) = \sum_{v \in \mathcal{V}} p_{e}(x_i=v|\pmb{x}_{<i}) \log p_{e}(x_i=v|\pmb{x}_{<i})
\end{equation}

We explore two techniques for detecting patch boundaries based on entropies $H(x_i)$. The initial approach selects points whose entropy exceeds a specified global threshold, as depicted in~\autoref{fig:patching}. The alternative method locates points where entropy is significantly higher compared to the immediately preceding value. This latter strategy may also be viewed as pinpointing points that interrupt an otherwise approximately monotonic decline in entropy within the patch.

\begin{align*}
    \text{Global Constraint}\hspace{1cm}H(x_t)& > \theta_g\\
    \text{Approx. Monotonic Constraint}\hspace{1cm}H(x_t)& - H(x_{t-1}) > \theta_r
\end{align*}

Patch boundaries are determined through a simple preprocessing procedure that occurs as part of the data loading phase. This approach contrasts with the method used by~\citet{nawrot-etal-2023-efficient}, which involves training a classifier to detect patch boundaries based on entropy calculations. In our experiments~(\S\ref{section:experiments}), we conduct a comparison of these two strategies for differentiating between bytes of low and high entropy.

\subsection{The Byte-Pair Encoding (BPE) Tokenizer and Incremental Patching}
\label{section:incremental-patching}
\label{section:bpe-patch}

Numerous contemporary llms, such as our baseline Llama 3, employ subword tokenization methods like bpe~\citep{gage1994new, sennrich-etal-2016-neural}. Throughout this work, we define ``tokens'' as collections of bytes mapped from a \textit{finite} vocabulary established before training, in contrast to ``patches,'' which are on-the-fly aggregations of sequences that do not rely on a predefined vocabulary. A fundamental distinction lies in the fact that, when using tokens, the model cannot directly access the raw byte-level features.

A key advancement of {\textsc BLT} relative to tokenization-based models lies in its reconfiguration of the balance between vocabulary size and computational requirements. In conventional llms, expanding the vocabulary typically results in tokens representing larger textual spans, which reduces the number of processing steps but simultaneously inflates the dimensionality of the model’s output layer. This compromise restricts tokenization-based strategies from achieving substantial variation in both token granularity and inference expense. For instance, Llama 3 achieves an increase in average token size—from 3.7 to 4.4 bytes—by expanding its embedding table to four times the size of that used in Llama 2.

During generation, {\textsc BLT} must determine at each position in the byte sequence whether a patch boundary has been reached, as this affects whether the Latent Transformer will be called for additional computation. This choice must be made solely based on information available up to the current point, without knowledge of bytes that have not yet been produced. Consequently, the patching mechanism cannot utilize future context when deciding how to divide the byte sequence. More formally, a patching scheme $f_p$ is said to support incremental patching if it adheres to the following condition:
$$f_p(\pmb{x}_{<i}) = f_p(\pmb{x})_{<i}$$
bpe does not qualify as an incremental patching scheme, since the tokenization of a prefix can vary depending on what follows in the sequence, thereby failing to meet this criterion\footnote{Appending a dedicated delimiter token to mark patch boundaries could make bpe incremental; however, this approach extends the byte-sequence length.}.

\section{{\textsc BLT} Architecture}
\label{section:architecture}
The {\textsc BLT} system is structured with a large, global autoregressive language model responsible for handling patch-level representations. In addition, it incorporates two lightweight local models tasked with mapping byte sequences to patches and reconstructing patch representations back into byte sequences, as shown in~(\autoref{fig:arch_high_level}).

\subsection{Latent Global Transformer Model}

\textit{The Latent Global Transformer} refers to an autoregressive transformer architecture, denoted as $\mathcal{G}$, consisting of $l_{\mathcal{G}}$ layers. This model receives a series of latent input patch embeddings, $p_j$, and transforms them into a corresponding sequence of output patch embeddings, $o_j$. In this work, indices $j$ are consistently used for patches, while $i$ indexes bytes. The architecture employs a block-causal attention mask~\citep{dubey2024llama}, which ensures that each patch in the current document attends only to itself and to previous patches up to the present one. Notably, the global model accounts for most of the computational cost—both during pre-training and inference. Consequently, by deciding under which circumstances to apply this component, we are able to manage and adapt computational resource allocation across various segments of the input and output, depending on the underlying complexity.

\subsection{Local Encoder}
\label{section:local}

\textit{The Local Encoder Model}, represented as $\mathcal{E}$, comprises a compact transformer architecture containing $l_{\mathcal{E}} << l_{\mathcal{G}}$ layers. Its central function is to rapidly project an input sequence of bytes $b_i$ into rich patch representations $p_j$. Distinct from the standard transformer, this model incorporates a cross-attention mechanism following each transformer layer, enabling pooling of byte-level encodings into patch-level representations~(\autoref{fig:crossattn}). 

Initially, input bytes $b_i$ are projected into vectors $x_i$ using an embedding matrix of dimension $\mathbb{R}^{256 \times h_{\mathcal{E}}}$. These embeddings may be further supplemented with auxiliary hash-embeddings~(\S\ref{sec:hash-emb}). Subsequent layers—comprising alternating transformer and cross-attention modules~(\S\ref{sec:enc-cross-attn})—sequentially transform the embeddings into patch representations $p_i$, which are ultimately consumed by the global transformer $\mathcal{G}$.

Transformer layers employ a \textit{local block causal} attention masking scheme, restricting each byte’s attention to the $w_{\mathcal{E}}$ immediate preceding bytes. This window of attention may span across varying patch boundaries, though it is confined within the same document. The ensuing sections elaborate on the embedding process and the specifics of the cross-attention component.

\subsubsection{Encoder Hash n-gram Embeddings}
\label{sec:ngram-emb}
\label{sec:hash-emb}
Capturing informative and resilient representations at every position $i$ necessitates the integration of context derived from prior bytes. In {\textsc BLT}, this is accomplished by representing not only the individual byte $b_i$, but also its participation within a sequence, specifically a byte n-gram. At each position $i$, our approach begins by generating byte-grams

\begin{align}
    g_{i,n}=\{b_{i-n + 1},\ldots, b_i\}
\end{align}

for every byte position $i$ and for values of $n$ ranging from three up to eight.\footnote{
Byte-grams with length $n$ or more are not considered if $i<n$. }

Next, we define hash-based $n$-gram embeddings: each byte $n$-gram is mapped by a hash function to an index in an embedding matrix $E_{n}^{hash}$ with fixed dimensionality, separately for each $n \in\{3,4,5,6,7,8\}$~\citep{bai2010learning}. The corresponding embedding vector is added to the original byte embedding, and the combined representation is normalized before being input to the local encoder architecture. The enhanced embedding is computed as:

\begin{align}
e_i &= x_i + \sum_{n = 3,...,8}E_{n}^{hash}(\text{Hash}(g_{i,n})) \\
\text{where, Hash}(g_{i,n}) &= \text{RollPolyHash}(g_{i,n}) \% |E_{n}^{hash}|
\end{align}

The value $e_i$ is normalized by dividing by the sum of the number of $n$-gram sizes and one. The RollPolyHash function, as specified in Appendix~\ref{appendix:rollpolyhash}, is employed in this context. Section~\ref{section:ablations} investigates the impact of varying the parameters of $n$-gram hash embeddings, specifically the choice of $n$ and the size of the embedding table, on flop-controlled scaling law behaviors. Beyond hash-based $n$-gram embeddings, we also explored $n$-gram embeddings based on frequency information; further information about this approach is provided in Appendix~\ref{appendix:ngrams}.

\subsubsection{Encoder Multi-Headed Cross-Attention}
\label{sec:enc-cross-attn}
Our approach largely mirrors the input cross-attention component found in the Perceiver architecture~\citep{jaegle2021perceiver}; however, the primary distinction lies in the use of variable patch representations as the latent variables instead of a static latent set~(\autoref{fig:crossattn}). Each patch only attends to the bytes associated with that particular patch. For every patch $p_j$, the module constructs a query vector by aggregating the byte representations assigned to $p_j$—commonly via a pooling operation—after which a linear projection, $\mathcal{E}_{C} \in \mathbb{R}^{h_{\mathcal{E}} \times (h_{\mathcal{E}} \times U_{\mathcal{E}})}$, is applied, where $U_{\mathcal{E}}$ is the total count of encoder cross-attention heads. Letting $f_{\text{bytes}}(p_j)$ represent the ordered sequence of bytes for patch $p_j$, we compute

\begin{align}
P_{0,j} &= \mathcal{E}_{C}(f_\text{bytes}((p_j)), f~ \text{is a pooling function} \\
P_l &= P_{l-1} + W_o\left(\text{softmax}\left(\frac{QK^T}{\sqrt{d_k}}\right)V\right) \\
\text{where } Q_j &= W_q(P_{l-1,j}), K_i = W_k(h_{l-1,i}), V_i = W_v(h_{l-1,i}) \\
h_l &= \text{Encoder-Transformer-Layer}_l(h_{l-1})
\end{align}

In this context, $P \in \mathbb{R}^{n_p \times h_{\mathcal{G}}}$ denotes the $n_p$ patch representations designated for processing by the global model. These patch representations are initially formed by aggregating the byte embeddings $e_i$ associated with each individual patch $p_j$. The matrices $W_q$, $W_k$, $W_v$, and $W_o$ serve as the projection weights for queries, keys, values, and outputs, respectively. Here, keys and values are generated as linear projections of the byte representations $h_i$ output by the previous network layer (or $e_i$ at the first layer). To ensure each query $Q_j$ only accesses information within its corresponding patch, a patch-aware masking method is applied so that each query attends exclusively to the keys and values of its own patch. Given that we implement multi-headed attention on $Q, K$, and $V$, and that patch representations typically have a larger dimension ($h_{\mathcal{G}}$) relative to $h_{\mathcal{E}}$, $P_l$ is structured to maintain multiple attention heads of size $h_{\mathcal{E}}$ throughout the cross-attention process, and these are subsequently concatenated to achieve the full dimensionality $h_{\mathcal{G}}$. Furthermore, we precede the attention with LayerNorm on the queries, keys, and values, omitting any positional embeddings in this module. The architecture concludes with a residual connection enveloping the cross-attention block.

\subsection{Local Decoder} 

The local decoder $\mathcal{D}$, akin to the local encoder in design, employs a compact transformer architecture with a number of layers $l_{\mathcal{D}}$ substantially smaller than $l_{\mathcal{G}}$. Its function is to map a sequence of global patch representations $o_j$ back into the original raw byte sequence $y_i$. The prediction of each byte within the sequence is conditioned on previously generated bytes and leverages the hidden states provided by the local encoder corresponding to those bytes. The architecture consists of $l_{\mathcal{D}}$ layers, each comprising a cross-attention block followed by a transformer block. Initially, the decoder’s cross-attention mechanism constructs byte-level representations from the given patch representations, which are then refined by subsequent transformer layers as they process the evolving byte sequence.

\subsubsection{Decoder Multi-headed Cross-Attention} 

Within the decoder's cross-attention mechanism, the queries are derived from the byte-representations, whereas the patch representations serve as the keys and values—a reversal of their conventional roles. For this cross-attention phase, the byte-representations are initialized using the byte embeddings obtained from the final encoder layer, namely $h_{l_{\mathcal{E}}}$. The byte-representations at each decoder layer $l$, denoted as $d_{l,i}$, are subsequently updated according to the following:

\begin{align}
D_0 &= h_{l_{\mathcal{E}}} \\
B_l &= D_{l-1} + W_o\left(\text{softmax}\left(\frac{QK^T}{\sqrt{d_k}}\right)V\right), \\
  &\text{where } Q_i = W_q(d_{l-1,i}), K_i = W_k(\mathcal{D}_C(o_j)), V_i = W_v(\mathcal{D}_C(o_j)) \\
  D_l &= \text{Decoder-Transformer-layer}_l(B_l)
\end{align}

Here, $W_k$ and $W_v$ denote the key and value projection matrices, respectively; these are applied to the final patch embeddings $o_j$ from the global model through a linear transformation followed by a split operation $\mathcal{D}_C$. The query projection matrix $W_q$ is applied to the byte-level representations $d_{l-1}$ output by the previous decoder transformer layer, or $h_{l_{\mathcal{E}}}$ for the initial layer. The matrix $W_o$ projects the outputs, resulting in $B \in \mathbb{R}^{h_{\mathcal{D}} \times n_b}$, with $n_b$ indicating the total number of output bytes. The subsequent set of decoder representations, $D_l$, is produced by passing $B$ through a decoder transformer layer. Consistent with the cross-attention mechanism used in the local encoder, this architecture uses multi-head attention, applies pre LayerNorm, omits positional embeddings, and incorporates a residual connection encircling the cross-attention mechanism.

\section{Experimental Setup}
\label{section:experiments}
We perform rigorously controlled experiments to directly compare {\textsc BLT} and tokenization-based models, taking care to ensure that {\textsc BLT} does not benefit from potential sequence length advantages.

\subsection{Pre-training Datasets} In our experiments, all model variants are initially trained on two primary datasets: (1) The Llama 2 dataset~\citep{touvron2023llama}, containing 2 trillion tokens sourced from diverse, publicly available materials. This data undergoes extensive cleaning and filtering processes to enhance quality. (2) {\textsc BLT}-1T: A novel corpus consisting of 1 trillion tokens, compiled from a range of public sources and incorporating a portion of the pre-training data from Datacomp-LM~\citep{li2024datacomp}. The Llama 2 dataset serves as the foundation for scaling law analyses regarding the optimal number of training tokens, following the methodology of~\cite{dubey2024llama}, and thus informs architectural optimizations for {\textsc BLT}. In contrast, {\textsc BLT}-1T is used for end-to-end pre-training, enabling head-to-head downstream performance comparisons with Llama 3. It is important to note that both datasets exclude any information derived from Meta products or services. For our baseline evaluations involving tokenizer-based models, the Llama 3 tokenizer with a 128K vocabulary size is employed; this configuration yields improved baseline performance in our experiments relative to the Llama 2 tokenizer.

\subsection{Entropy Model} For our experiments, the entropy model employed is a byte-level language model, trained using the identical training data as the comprehensive {\textsc BLT} model. By default, this model configuration consists of a 100M-parameter transformer comprising 14 layers, a hidden size of 512, and sliding window attention over 512 bytes. Other hyperparameters are aligned with those used in both our local and global transformer architectures. As explored in \autoref{section:ablations}, we varied model capacity, receptive field length, and model design. Notably, in scenarios where the model's receptive field is sufficiently constrained, the resulting trained entropy model can be efficiently implemented via a lookup table.

\subsection{Entropy Threshold and Equalizing Context Length} For models employing entropy-driven patching, we determine a patching threshold so as to produce a specified average \textit{patch size} on the pretraining data mixture. In the context of {\textsc BLT}, as opposed to token-based approaches, the \textit{patch size} can be flexibly set, substantially affecting the model’s effective context window. To standardize the average context length and prevent potential biases favoring models with larger patch sizes, we keep the expected total number of bytes per batch fixed. As a result, models operating with larger patch sizes have their sequence lengths proportionally decreased. Specifically, we adopt an 8k byte context length for Llama 2 data and a 16k byte average context for the {\textsc BLT}-1T dataset, all while ensuring an average batch size of 16M bytes is maintained.

Although the mean batch size remains unchanged, dynamic patching techniques result in varying proportions of bytes per patch during data loading. To optimize efficiency, our {\textsc BLT} training system groups batches based on patches, thereby circumventing the need for padding within the computationally intensive latent transformer. This strategy guarantees a uniform patch count across all batches. To manage memory usage and prevent unexpected spikes due to exceptionally large patches, we pad—potentially also truncate—input byte sequences to fixed lengths of 12k bytes for the Llama 2 dataset and 24k bytes for the {\textsc BLT}-1T dataset during training.

\subsection{Entropy Model Context}
\label{section:entropy-context}
Our empirical analysis reveals that implementing entropy patching tends to generate increasingly larger patches in tasks characterized by structured formats, such as multiple choice questions (refer to the MMLU patching example in \autoref{fig:mmlu_patching}). This tendency can be attributed to the lower entropy observed in repeated sequences within the entropy model context, which leads to more pronounced variations. Accordingly, for the large-scale {\textsc BLT}-Entropy experiment with a patch size of 4.5, we reset the entropy context by introducing new lines and employ the approximate monotonicity constraint, since this approach demonstrates greater resilience to “entropy drift” that arises due to fluctuating context lengths. It is important to note that this modification pertains solely to the entropy computation process, while our method for determining the entropy threshold remains unchanged.

\subsection{FLOPs Estimation} 
\label{section:flops}

Our approach to calculating transformer flops is primarily based on the methodology outlined in Chinchilla~\citep{hoffmann2022training}, which accounts for flops in the feed-forward components, the qkvo projections within self-attention mechanisms, as well as those arising from attention computation and the output projection. One important distinction in our analysis is the assumption that the input embedding layer functions via a highly efficient lookup, as opposed to employing a dense matrix multiplication, which effectively renders it a 0-flop process. In line with earlier studies, we approximate that the backward pass consumes twice as many flops as the forward pass.

To determine the flops \textit{per byte} metric for {\textsc BLT} architectures, we aggregate the floating-point operations required by each component: the local encoder transformer, the global latent transformer, and the local decoder transformer, as well as the contributions from the cross-attention modules found within both encoder and decoder stages. This can be formalized as follows:

\begin{align}
    \text{FL}_{\text{{\textsc BLT}}} &= \frac{1}{n_p}\, \text{Transf. FL}(h_{\mathcal{G}}, l_{\mathcal{G}}, m=n_{ctx}/n_p,V=0) \\
   &+ \text{Transf. FL}(h_{\mathcal{E}}, l_{\mathcal{E}}, m=w_{\mathcal{E}}, V=0) \\
   &+ \text{Transf. FL}(h_{\mathcal{D}}, l_{\mathcal{D}}, m=w_{\mathcal{D}}, V=256) \\ 
    &+ k/n_p \times \text{Cross Attn. FL}(h_{\mathcal{E}}, l_{\mathcal{E}}, m=n_p, r=n_p/k) \\ 
   &+ \text{Cross Attn. FL}(h_{\mathcal{D}}, l_{\mathcal{D}}, m=k, r=k/n_p) 
\end{align}

In this context, $n_{ctx}$ denotes the sequence length measured in bytes, $n_p$ signifies the patch size, and $r$ specifies the proportion of queries relative to keys/values. The parameter $k$ reflects the quotient between the patch dimension and the byte dimension, representing how many local model partitions are combined to yield a comprehensive model representation (e.g., $k=2$ as depicted in \autoref{fig:crossattn}). The variable $V$ refers to the size of the vocabulary associated with the output projection, which is utilized exclusively in the local decoder. The value of $m$, which represents the context length adopted in the attention mechanism, differs depending on whether the module operates on the byte-level or the patch-level sequence. Consequently, the computation of attention-related flops is adjusted for each component. Detailed formulations for calculating the Transformer-FLOPs and Cross-Attention FLOPs can be found in Appendix~\ref{section:flops-appendix}.

\subsection{Bits-Per-Byte Estimation} The notion of perplexity is inherently tied to a specific tokenizer, as it quantifies the model’s uncertainty at the token level. Therefore, to enable fair comparisons between models that operate at the byte level and those that use tokens, we adopt the Bits-Per-Byte (BPB) metric, in line with earlier research~\citep{xue2022byt5, yu2023megabyte, wang2024mambabyte}. BPB serves as a version of perplexity that does not depend on any tokenizer. It is defined as follows:

\begin{align}
    \text{BPB}(x)&= \frac{\mathcal{L}_{CE}(\pmb{x})}{\ln(2)\cdot n_{\text{bytes}}}
\end{align}

In this formulation, the total cross-entropy loss computed over the dataset $\pmb{x}$ is divided by the number of bytes in $\pmb{x}$, and further scaled by a constant. This yields a standardized, tokenizer-agnostic measure of uncertainty.

\subsection{Transformer Architecture Hyperparameters} 

In configuring all transformer blocks within {\textsc BLT}—encompassing both local and global models—we adhere closely to the architectural choices found in Llama~3~\citep{dubey2024llama}. Specifically, feed-forward layers employ the SwiGLU activation function~\citep{swiglu}, while rotary positional embeddings (RoPE)~\citep{RoPe} with $\theta=500000$~\citep{xiong2024effective} are incorporated solely within self-attention mechanisms. Layer normalization is carried out using RMSNorm~\citep{RMSNorm}. For self-attention layers utilizing standard fixed attention masks, such as \textit{block causal} or \textit{fixed-window block causal}, Flash Attention~\citep{dao2022flashattention} is implemented, with a window size of 512 applied to all fixed-width attention masks. In contrast, for cross-attention layers that require dynamically computed patch-dependent masks, Flex Attention\footnote{\url{https://pytorch.org/blog/flexattention}} is adopted, providing fused kernels that substantially enhance training efficiency.

\subsection{{\textsc BLT}-Specific Hyperparameters}

To evaluate the performance of {\textsc BLT} models, our experiments explore two primary areas: scaling laws and downstream task performance, analyzing models with varying sizes: 400M, 1B, 2B, 4B, and 8B parameters. Full details of the architectural hyperparameters are provided in Appendix~Table~\ref{tab:arch}.

For all {\textsc BLT} models, we employ max-pooling for initializing queries in the first cross-attention layer of the local encoder. We configure hashing using $500,000$ hashes and a single hash function, selecting n-gram sizes between 3 and 8. The learning rate is consistently set at $4e-4$ across all models. This learning rate selection is motivated by a sweep at the 400M and 1B model sizes over the $1e-3$ to $1e-4$ range, which revealed identical optimal rates for both the token-based and {\textsc BLT} models. In scaling experiments with Llama-2 datasets, we set the batch sizes according to the guidelines in~\cite{dubey2024llama}, or ensure the byte-size equivalent.

Optimization throughout utilizes AdamW~\citep{adamw}, where $\beta_1$ and $\beta_2$ are set to 0.9 and 0.95, respectively, and $\epsilon=10^{-8}$. We schedule a linear warm-up phase lasting 2000 steps, followed by a cosine decay of the learning rate to zero. The models further employ a weight decay of 0.1, with global gradient norm clipping enforced at 1.0.

\section{Scaling Trends}
\label{section:scaling}

We provide a comprehensive analysis of the scaling behavior of byte-level models, offering insights that can guide the further development of {\textsc BLT} models. Our scaling investigation seeks to overcome the constraints seen in earlier byte-level model studies by: (a) examining trends under compute-optimal training conditions, (b) training equivalently sized 8B models with substantial training sets (reaching up to 1T tokens/4T bytes) and assessing their performance on downstream tasks, and (c) quantifying scaling characteristics when inference costs are held constant. Subsequent sections will delve into distinct benefits associated with modeling byte sequences.

\subsection{Parameter Matched Compute Optimal Scaling Trends}
\label{section:parameter-matched-scaling}
Leveraging the Llama 2 dataset, we conduct training of both \textit{compute-optimal} bpe models and their {\textsc BLT} counterparts at four distinct scales, spanning 1B to 8B parameters. To assess scaling behaviors, we compare the number of training flops with language modeling accuracy on a representative portion of the training mixture. The bpe models utilize the parameter-to-data ratio specified as optimal by Llama~3~\citep{dubey2024llama}. This \textit{compute-optimal} regime is constructed to yield maximal performance on the training distribution for a fixed computation budget~\citep{hoffmann2022training}, serving as a rigorous baseline for evaluation. Alongside each bpe baseline, we train a {\textsc BLT} model with identical data, employing a Latent Transformer whose size and architecture are aligned with those of the corresponding bpe Transformer.

As shown in~\autoref{fig:scaling-compute-opt} (right), {\textsc BLT} models consistently equal or surpass the performance of their bpe-based equivalents, a pattern that persists across increases in model scale and computational resources. To our knowledge, {\textsc BLT} represents the first byte-level Transformer architecture to exhibit scaling characteristics comparable to those of BPE-based models under compute-optimal conditions. This finding substantiates our hypothesis that the parameter-to-compute ratio deemed optimal for bpe-based models is also suitable for {\textsc BLT}, or at minimum, does not significantly deviate from it.

Advancements in model architecture and the adoption of dynamic patching are both essential for aligning with bpe scaling patterns. In ~\autoref{fig:scaling-compute-opt} (left), we present a comparison between models based on space-patching methods and Llama 3. Our approximation of SpaceByte \citep{slagle2024spacebyte} employs {\textsc BLT} space-patching while omitting n-gram embeddings and cross-attention components. While SpaceByte exhibits performance gains relative to Megabyte, it still lags significantly behind Llama 3. On the right side of \autoref{fig:scaling-compute-opt}, we display the gains that arise from incorporating both novel architectural features and dynamic patching. At larger scales, {\textsc BLT} models achieve results that are competitive with leading tokenizer-based models like Llama 3.

Additionally, we find that the selection of tokenizer has a notable impact on the effectiveness of tokenizer-based models. Specifically, models utilizing the Llama-3 tokenizer demonstrate superior performance compared to those that use the Llama-2 tokenizer on identical datasets.

In summary, our {\textsc BLT} model exhibits intermediate scaling behavior between Llama 2 and Llama 3, particularly when much larger patch sizes are employed. While the average token sizes for the bpe tokenizers in Llama 2 and 3 are 3.7 and 4.4 bytes, respectively, {\textsc BLT} achieves comparable scaling dynamics with mean patch sizes of 6 or even 8 bytes. Since inference flops are inversely related to the chosen patch size, opting for an 8-byte patch can result in up to a 50\% reduction in inference flops. Furthermore, increasing the patch size appears to benefit model performance as both model and dataset scale up. Although {\textsc BLT} with an 8-byte patch size initially underperforms relative to the bpe-based Llama 2 at the 1B parameter level, it surpasses bpe performance by the 7B parameter scale. These findings indicate that larger patch sizes may yield superior outcomes at higher model scales and suggest that scaling to even greater patch sizes may become viable with further increases in model capacity and available training compute.

\subsection{Beyond Compute Optimal Task Evaluations}
\label{section:evals}
To further investigate scaling behavior, we continue training an 8B {\textsc BLT} model past the point of compute-optimal scaling, utilizing the {\textsc BLT}-1T dataset, which is both larger and of higher quality. We then assess the model's performance on a broad range of standard benchmarks for classification and generative abilities. The following tasks are chosen to evaluate common sense reasoning, knowledge, and code synthesis:
\paragraph{Classification tasks} assessed include ARC-Easy (0-shot)~\citep{Eval_ARC}, Arc-Challenge (0-shot)~\citep{Eval_ARC}, HellaSwag (0-shot)~\citep{Eval_hellaswag}, PIQA (0-shot)~\citep{Eval_piqa}, and MMLU (5-shot)~\citep{hendrycks2020measuring}. For these evaluations, we use a prompt-scoring approach that computes the likelihood of each answer choice, and we present average accuracy as the metric.
\paragraph{Coding related generation tasks:} To assess large language model capabilities in Python code generation, we report pass@1 metrics on MBPP (3-shot)~\citep{Eval_MBPP} and HumanEval (0-shot)~\citep{Eval_HumanEval}.

In \autoref{tab:evals}, we present a comparison among three models trained using the {\textsc BLT}-1T dataset: a Llama 3 model based on the bpe tokenizer,\footnote{We opt for the Llama 3 tokenizer with its larger 128k vocabulary, as it demonstrates improved performance relative to Llama 2’s 32k vocabulary.} along with two versions of the {\textsc BLT} model. One version adopts the space-patching approach ({\textsc BLT}-Space), while the other employs an entropy-based patching method ({\textsc BLT}-Entropy) that enforces an approximate monotonicity constraint and resets the context with newlines, following the procedure outlined in~\autoref{section:entropy-context}. All three models are trained with an equal allocation of computational flops. Notably, during inference with {\textsc BLT}-Entropy, we reduce the entropy threshold from 0.6 to 0.1, observing improved task performance at the expense of requiring more inference steps.

The {\textsc BLT}-Entropy model achieves superior results compared to the Llama 3 model in 4 of the 7 evaluated tasks, even though both models are trained on an equal volume of data. This enhanced performance is likely attributable to two primary factors: (1) more efficient utilization of training compute made possible by dynamic patching, and (2) the explicit representation of information at the byte level rather than at the token level.

Conversely, {\textsc BLT}-Space lags behind the Llama 3 tokenizer on all tasks except one. However, its increased average patch size of 6 bytes allows for a notable decrease in inference flops. By contrast, models utilizing bpe and entropy-patching approaches maintain an average patch size of about 4.5 bytes when evaluated on the training dataset mixture. Given an equivalent training budget, the model with the larger patch size is able to process 30\% more data than its counterparts. This expanded coverage could potentially result in {\textsc BLT} deviating further from the point of compute-optimality.

\subsection{Patches Scale Better Than Tokens} With {\textsc BLT} models, it is possible to scale both the model capacity and the patch size at the same time, without exceeding a fixed training or inference computational budget, and without increasing the dataset size. The flexibility to increase patch size arbitrarily is exclusive to patch-based architectures and allows these models to circumvent the efficiency limitations inherent to fixed-vocabulary token-based systems, as elaborated in Section~\ref{section:bpe-patch}. Employing larger patches reduces the required computation, permitting the allocation of saved resources toward expanding the global latent transformer, since its invocation frequency decreases.

We perform a fixed inference scaling investigation to evaluate whether, for a fixed inference budget, increasing model size while reducing the number of steps on larger patches yields superior performance compared to employing smaller models with more steps. Using initial model configurations with 400 million and 3.6 billion parameters paired with the Llama 2 tokenizer, we identify flop-matched models using both the Llama 3 tokenizer and {\textsc BLT}-Entropy architectures, employing average patch sizes of 6 and 8 bytes in the training datamix (refer to~\autoref{tab:fixed-inf-params} for comprehensive model specifications). For models utilizing an 8-byte patch size, the encoder stack comprises 3 layers in place of a single layer. Each model undergoes training across several allocated training flop budgets.

As illustrated in \autoref{fig:fixed_inference_scaling}, {\textsc BLT} architectures demonstrate superior scaling behavior compared to tokenization-based models across both inference flop classes. While BPE-based models tend to outperform at smaller training expenditures, {\textsc BLT} soon overtakes them just past the compute-optimal point. This suggests that allocating greater resources to pretraining—even if it is a one-time investment—can lead to a more capable model when working under a fixed inference computational budget. A notable case is the 8B parameter class, such as Llama 3.1, which is pretrained on data volumes nearly two orders of magnitude larger than what the compute-optimal threshold would suggest for its scale.

The point at which {\textsc BLT} surpasses token-based architectures has moved marginally closer to the compute-optimal threshold as model size increases in the higher flop category—shifting from 3x to 2.5x the compute-optimal allocation. In the same vein, the model employing a larger patch size of 8 demonstrates a sharper scaling trajectory within this high flop bracket, surpassing other variants at an earlier stage. As outlined in Section~\ref{section:parameter-matched-scaling}, scaling up patch sizes leads performance to align more closely with that of BPE models when using larger networks. This phenomenon can be partially explained by the reduced proportion of overall flops consumed by the byte-level Encoder and Decoder modules, which exhibit a slower scaling behavior compared to the Latent Transformer. For instance, when increasing the total number of parameters by a factor of 20—from 400M up to 8B—the local parameter count for {\textsc BLT} approximately doubles. This detail is noteworthy, as enlarging the patch size predominantly influences the flops attributable to the patch Latent Transformer, while leaving the byte-level components largely unaffected. Consequently, {\textsc BLT}-Entropy with ps=8 experienced an increase from 1.6x to 1.7x relative to the Llama 2 model size as model scale expanded.

In conclusion, our investigation into patch-length scaling reveals that the {\textsc BLT} patch-based architecture exhibits superior scaling properties when both the patch size and model capacity are increased together. Notably, these favorable scaling behaviors continue, and may even be enhanced, as the model size grows.

\section{Byte Modeling Improves Robustness} We further evaluate the robustness of {\textsc BLT} in comparison with token-based architectures that do not leverage explicit byte-level representations, and introduce a method for transforming pretrained token-based models to operate at the byte level.
\label{sec:robustness}

\subsection{Character-Level Tasks}

An initial impetus for developing byte-level models stemmed from their resilience to input noise at the byte level, as well as their ability to recognize and process the components within tokens—a challenge for many models dependent on tokenizers. To assess these attributes, we conducted supplementary evaluations using benchmarks designed to test both robustness against input noise and the capacity to recognize character-level information in English and various other languages, encompassing not only letters but also digits and phonemes. Table~\ref{tab:char_tasks} summarizes the findings from these experiments.

\paragraph{Noisy Data} To assess model robustness, we generate altered versions of the benchmark classification tasks from Section~\ref{section:evals} by introducing character-level noise. Our approach involves five specific noising techniques: (a) \textit{AntSpeak}, which converts the text into uppercase letters separated by spaces; (b) \textit{Drop}, where 10\% of the characters are randomly deleted; (c) \textit{RandomCase}, in which half the characters are randomly capitalized and the rest made lowercase; (d) \textit{Repeat}, where up to 20\% of the characters are repeated as many as four times; and (e) \textit{UpperCase}, which changes all text to uppercase. Each of these noise strategies is applied either to the prompt, the completion, or both—each considered a distinct evaluation task, with average performance reported. Table~\ref{tab:char_tasks} presents results on the noised version of HellaSwag~\citep{Eval_hellaswag}, revealing that {\textsc BLT} consistently exhibits superior robustness compared to tokenizer-based models. On average, {\textsc BLT} achieves an 8-point performance increase over the comparable model and even surpasses the Llama 3.1 model that was trained using a substantially larger corpus.

\paragraph{Phonology - Grapheme-to-Phoneme (G2P)} We evaluate how effectively {\textsc BLT} can convert a string of graphemes (letters encoding a word) into a corresponding sequence of phonemes that represent the word's spoken form. The outcomes for the G2P task in a 5-shot scenario, as detailed in Table \ref{tab:char_tasks} and utilizing Phonology Bench~\citep{suvarna2024phonologybench}, demonstrate that {\textsc BLT} achieves superior performance compared to the Llama 3 1T tokenizer-based baseline model.

\paragraph{CUTE}
To evaluate character-level comprehension, we test {\textsc BLT} on the CUTE benchmark~\citep{edman2024cute}, which features a range of tasks grouped into three principal areas: compositional understanding, orthographic similarity, and sequence manipulation. This suite of tasks is notably difficult for most models based on tokenization, which typically possess some awareness of how their tokens are spelled but encounter difficulties when required to exploit this knowledge for text manipulation. As demonstrated in Table~\ref{tab:char_tasks}, {\textsc BLT}-Entropy achieves performance exceeding that of both BPE Llama 3 models by over 25 points. The model exhibits particularly strong capabilities in tasks requiring character-level manipulation, reaching 99.9\% accuracy on both spelling-related tasks. Remarkably, these gains occur even though {\textsc BLT} was trained on only one-sixteenth the data used for Llama 3.1, suggesting that acquiring character-level representations remains a significant challenge for BPE models. Figure~\ref{fig:cute_examples} provides examples where the Llama 3 tokenizer model fails, yet the {\textsc BLT} model delivers correct outputs. The only areas where BPE-based models retain a marginal advantage are in word deletion and insertion; these forms of manipulation can pose unique challenges to byte-level models, though the performance difference is minimal. It may also be more tractable for a model to build word-level competencies from characters than the reverse. For all tasks, we adopt an identical evaluation framework and utilize Huggingface's original prompts. Additional prompt engineering could potentially enhance BPE model performance.

\paragraph{Low Resource Machine Translation} We assess the capabilities of {\textsc BLT} in both source-to-target and target-to-source translation for twenty-one low-resource languages across six major language families, each utilizing a range of scripts, as provided by the FLORES-101 benchmark~\citep{goyal2022flores}. SentencePiece BLEU scores are summarized in Table~\ref{tab:flores}. The findings indicate that {\textsc BLT} consistently surpasses a model employing the Llama 3 tokenizer, with a 2-point aggregate improvement for translations into English, and a 0.5-point gain when translating out of English. In translation tasks involving high-resource language pairs, {\textsc BLT} matches or modestly exceeds Llama 3’s performance. More notably, {\textsc BLT} demonstrates superior results over Llama 3 across many lower-resource language pairs, highlighting the benefit of byte-based modeling for handling long-tail byte sequence generalization.

\subsection{Training {\textsc BLT} from Llama 3}
\label{sec:distillation}
We investigate a methodology in which {\textsc BLT} models benefit from pre-trained tokenizer-based models to improve training speed and convergence. This is accomplished by initializing {\textsc BLT}'s global transformer parameters using those from a pre-trained Llama 3.1 model. Following this initialization, the global transformer weights are fine-tuned with a learning rate that is one-tenth of the rate applied to the local encoder and decoder components, using the optimal step count as determined for Llama 3. A comparison to a standard {\textsc BLT} baseline is provided in Table~\ref{tab:distill}. Results show that {\textsc BLT} initialized from Llama 3.1 markedly outperforms both the original Llama 3 and the baseline {\textsc BLT}, given the same compute budget (flops). Additionally, even when evaluated against the {\textsc BLT}-Entropy model—trained on a much larger dataset of 1T tokens, as reported in Table~\ref{tab:evals}—{\textsc BLT} from Llama 3.1 demonstrates better results on the MMLU task. This indicates that the proposed approach is highly effective in reducing training flops without compromising performance.

This approach can be interpreted as adapting tokenizer-dependent architectures into tokenizer-independent ones, essentially reconfiguring a pre-trained LLaMA 3.1 model into a {\textsc BLT} model. For thorough comparison, Table \ref{tab:distill} presents results from the original LLaMA 3.1, which was trained on 15T tokens, alongside the corresponding {\textsc BLT} generated from LLaMA 3. While our model shows only a minor reduction in accuracy on MMLU and HumanEval, it suffers a more pronounced decrease in performance across the remaining benchmarks. These results indicate that additional research is necessary to fully exploit the pre-trained model’s capabilities, especially with regard to optimizing the composition of training data and tuning other hyperparameters.

\section{Ablations and Discussion}
\label{section:ablations}
Here, we analyze ablations that inform the architectural selections behind {\textsc BLT} as well as decisions regarding patching strategies and hyper-parameter choices for the {\textsc BLT} model with 8B parameters trained on the {\textsc BLT}-1T dataset.

\paragraph{Entropy Model Hyper-parameters} To evaluate how the size of the entropy model and length of the context window influence scaling behavior, we experiment with byte-level entropy transformer models ranging in size from 1 million to 100 million parameters and utilize context windows between 64 and 512 tokens. We visualize the relationship between bpb and the scaling of training FLOPs using curves derived from our $400m$ and $1b$ {\textsc BLT} models trained on the Llama-2 dataset, as shown in \autoref{fig:entropy_ablation}. Our observations indicate that both increasing entropy model size and extending context window length enhance scaling performance, although the improvements taper off past approximately 50 million parameters.

\paragraph{Types of Patching}

We conduct an ablation study on the four patching strategies described in Section~\ref{section:patching}: 1) Strided Patching utilizing strides of 4 and 6, 2) Patching triggered by whitespace, 3) Patch selection guided by BPE Tokenizer using the Llama 3 tokenizer, and 4) Entropy-driven patching informed by a compact byte-level language model.

Although dynamic patching shortens the sequence length, we regulate this length to ensure that all patching strategies have comparable context lengths. To avoid confounding influences that arise from modeling larger contexts, every model is exposed to the same expected number of bytes per sequence during both training and inference. The outcomes of these ablation studies are illustrated in Figure~\ref{fig:scaling-compute-opt}. Among the evaluated approaches, all forms of patching except static patching show superior performance, with space patching achieving results nearly equivalent to those of dynamic entropy based patching.

In \autoref{tab:evals_ablation}, we provide comparative benchmark results for {\textsc BLT} models that utilize tokenizer-based approaches, space patching, and entropy-based patching, all trained on the Llama 2 dataset with the optimal training duration as suggested by~\citet{dubey2024llama}. While space patching represents a more straightforward approach that circumvents the need to execute an entropy model dynamically during training, our findings indicate that the improvements linked to entropy-based patching in the context of scaling behavior (see Section~\ref{section:scaling}) persist and are reflected in downstream evaluation tasks as well.\footnote{Results for space patching are derived from preliminary experiments lacking cross-attention, but analogous patterns are also evident when cross-attention is incorporated.}

\paragraph{Cross-Attention}
\autoref{tab:crossattn_ablations_1b} presents an ablation study on the insertion of cross-attention mechanisms at different locations within the encoder and decoder of {\textsc BLT}. Specifically, for the encoder-side cross-attention, we evaluate three methods for initializing the queries: (1) utilizing an identical learned embedding for each global state, (2) applying a hash embedding derived from the patch's byte values, and (3) employing a pooling operation over the encoder’s hidden representations corresponding to the patch bytes at the specified encoder layer.

Our results indicate that implementing cross-attention within the \textit{decoder} yields the greatest effectiveness. While the encoder experiences minor gains from cross-attention, these benefits are observed exclusively when queries are initialized through pooling. Furthermore, our analysis reveals that cross-attention provides notable advantages on Common-Crawl data, with the effect becoming more pronounced as patch size increases.

\paragraph{n-gram Hash Embeddings} 

We conduct an ablation study using n-gram hash embedding vocabulary sizes of 0, 100K, 200K, and 400K, with results summarized in Table~\ref{tab:ngram_ablations_1b}. The addition of hash embeddings consistently yields improvements across all evaluated domains, with the most notable gains observed on Wikipedia and Github datasets (a 0.04 bpb increase, compared to only 0.01 bpb after 15k steps at 8B parameters). At the 8B scale, reducing the hash size from 500K to 300K led to an insignificant change in performance, approximately 0.001 bpb after 15k training steps. These findings suggest that hash embeddings are essential for enabling {\textsc BLT} to reach parity with tokenizer-based model performance; nonetheless, improvements plateau beyond 300K hashes. Furthermore, the observed benefits from hash embeddings seem largely independent of those obtained through cross-attention, with both techniques enhancing results on distinct datasets.

\paragraph{Local Model Hyperparameters} Table \ref{tab:local_layers} presents an ablation study examining different configurations for the number of layers within the local encoder and decoder. Results indicate that, when utilizing hash n-gram embeddings, {\textsc BLT} achieves strong performance with a minimal encoder architecture—specifically, a single layer—while benefiting from a more substantial decoder.

\section{Related Work}
\label{section:related_work}

\textbf{Character-Level RNNs:} Since the advent of neural architectures, Character Language Modeling has remained a key area of interest~\citep{sutskever2011generating,mikolov2012subword,graves2013generating}, particularly due to its ability to naturally handle out-of-vocabulary terms, eliminating the need for back-off techniques. \cite{kim2016character} introduced a model that operates exclusively on input-side characters, employing convolutional and highway layers followed by LSTM-based RNNs. This approach demonstrated competitive performance against the leading RNN-based language models for English at the time, and it surpassed them in morphologically rich languages—highlighting a major benefit of character-level LLMs. Similarly, \cite{kenter2018byte} developed byte-level LSTM architectures for machine comprehension, which showed superiority over word-level systems on morphologically complex languages such as Turkish and Russian. In a comparable vein, \cite{zhang2015character} leveraged convolutional networks on character inputs for classification, surpassing word-level baselines in certain scenarios. To advance character-level modeling further, \citet{chung2019hierarchical} proposed hierarchical LSTM models with integrated boundary detectors at multiple layers, allowing the model to uncover latent hierarchical text structures and yielding performance improvements. Diverging from attention-based methods, ByteNet~\cite{kalchbrenner2016neural} approached machine translation with convolutional architectures at the character level.

\textbf{Character-Level Transformers:} The introduction of attention mechanisms in transformer models~\citep{vaswani2017attention} alongside subword tokenization strategies~\citep{sennrich-etal-2016-neural} brought marked advancements to neural language modeling and a variety of benchmark tasks. Nevertheless, the reliance on word and subword units imposes a certain inductive bias regarding the representational granularity at which these models operate. Seeking to bridge the gap between transformer architectures and earlier successes with character-based language models, \citet{al2019character} explored training highly deep transformer networks with the integration of auxiliary losses. Their work demonstrated that such transformer-based character-level models surpassed previous LSTM-based counterparts, yet they remained notably behind their word-level LLM counterparts in terms of performance. Similarly, GPT-2~\citep{radfordgpt2} reported that, when evaluated on expansive datasets such as the 1 billion word benchmark, byte-level language models failed to achieve the competitive performance observed in word-level models.

Although \cite{Choe2019BridgingTG} showed that byte-level LLMs utilizing transformers can surpass subword-level LLMs with similar parameter counts, these models are much more computationally intensive and require significantly longer training times. In a related effort, \cite{el2020characterbert} trained a BERT-based model, CharFormer, which generates word representations through convolutions over character embeddings. This method leads to performance gains in the medical domain, but also incurs considerable computational overhead. Meanwhile, \cite{clark2022canine} introduced CANINE, an encoder-only model with 150M parameters that processes raw character sequences. CANINE employs a deep transformer architecture akin to our global model, along with a combination of local transformers and strided convolutions to reduce the character sequence length, outperforming the comparable token-level model (mBERT) on various multilingual downstream benchmarks. In ByT5~\citep{xue2022byt5}, the authors experimented with byte-level encoder-decoder models that completely avoided patching mechanisms. While these models demonstrated greater resilience to noisy inputs and remained competitive with token-based models using one quarter of the data, the omission of patching resulted in computationally expensive attention calculations over each byte, making the approach resource-intensive. The direct modeling of bytes instead of subwords increases input sequence lengths, which poses challenges for efficiently scaling byte-level models. More recently, taking advantage of the Mamba Architecture~\citep{gu2023mamba}—which enables maintaining a fixed-size memory state across long contexts—\cite{wang2024mambabyte} developed a byte-level Mamba model that also eschews patching, outperforming comparable byte-level transformers in flop-controlled experiments at the 350M parameter scale, as measured by bits-per-byte across multiple datasets.

\textbf{Patching-based approaches:} Leveraging patching techniques has shown promise in reducing the otherwise high computational demands of byte-level LLMs, all while maintaining competitive performance levels. Early research has highlighted success stories at relatively small model scales and limited data volumes. For instance, \cite{nawrot-etal-2022-hierarchical} introduced the hourglass transformer, which incorporates static patching with both downsampling and upsampling strategies, resulting in superior performance over alternative byte-level models at the 150M parameter scale. Building on this, \cite{nawrot-etal-2023-efficient} introduced dynamic patching methods, such as a learned end-to-end boundary predictor, another version trained with tokenizer supervision, and an entropy-based approach akin to {\textsc BLT}. Their results demonstrate improvements over standard transformer models from that period on language modeling benchmarks, particularly at a 40M parameter scale using approximately 400M training tokens. Additionally, \cite{lester2024training} explored sequence compression through arithmetic coding, achieving greater compression ratios than BPE and employing an equal-info windows methodology. This enabled them to surpass byte-level baselines in language modeling tasks, though they still trailed behind models using subword tokenization.

Our research is primarily motivated by MegaByte~\citep{yu2023megabyte}, a causal decoder-only large language model that employs a fixed approach to patching by statically segmenting and concatenating representations in order to translate between bytes and patches. MegaByte utilizes a local module within the decoder to revert from patches back to the byte level. The original authors demonstrate that MegaByte achieves performance comparable to tokenizer-based models when using 1B parameters on datasets comprising approximately 400B bytes. In our empirical studies, we thoroughly evaluate MegaByte and observe that the reliance on static patching leads to suboptimal performance when compared to the most recent, compute-efficient tokenizer-based models under equal computational budgets. Our results show how {\textsc BLT} closes this performance gap. Similarly, \cite{slagle2024spacebyte} note these same limitations with MegaByte and propose enhancements by introducing patching strategies that exploit whitespaces and analogous space-like bytes, as well as incorporating a local encoder component. Their findings indicate improvements relative to transformer models using tokenization, particularly on domains like Github and arXiv at the 1B parameter scale. Our experiments using this method confirm these gains, but also highlight that further innovations in model architecture will be necessary to continue advancing byte-level model performance and to approach the capabilities of leading token-based architectures such as Llama 3.

\section{Limitations and Future Work}

In this study, our architectural decisions are guided by training models for the optimal number of steps, as established for Llama~3~\citep{dubey2024llama}. Nonetheless, it is important to note that these scaling laws were originally derived for BPE-level transformers, which may result in non-ideal ratios of data to parameter counts when applied to {\textsc BLT}. Determining the scaling laws specifically for {\textsc BLT} remains an open direction for future research, and doing so may reveal more advantageous scaling behaviors for our model. Furthermore, our experimental investigations are primarily limited to models with up to 1B parameters. It is plausible that as the parameter count increases to 8B and higher, the best architectural configurations may also shift, potentially enabling even greater performance at larger scales.

Current transformer codebases and libraries are primarily optimized for models that utilize tokenization-based transformer structures. Although we provide theoretically equivalent flop experiments and incorporate some efficient solutions—such as FlexAttention—to accommodate non-standard transformer layers, our implementations might still fall short of the runtime performance observed in tokenizer-based models and could require additional optimization to bridge this gap.

Although {\textsc BLT} employs an independently trained entropy model for patching, exploring end-to-end learning of the patching mechanism presents a promising avenue for future research. In Section \ref{sec:distillation}, we report preliminary results that suggest the feasibility of adapting tokenizer-based models like Llama 3—trained on over 10T tokens—to bytes. This is achieved by initializing and keeping the global transformer weights fixed. Further investigations in this area could yield approaches that preserve the advantages of bytefying while potentially surpassing the original tokenizer-based models' performance, all without necessitating full retraining.

\section{Conclusion}
\label{section:conclusion}

In this work, we introduce the Byte Latent Transformer (\textbf{BLT}), an innovative architecture that challenges the standard reliance on predetermined tokenization schemes in large language models. BLT incorporates a flexible, trainable mechanism to aggregate bytes into variable-length patches, thereby dynamically distributing computational effort according to the intrinsic complexity of the input data. This approach yields substantial gains in both computational efficiency and model resilience. Through a comprehensive scaling analysis, we show that {\textsc BLT} achieves parity with tokenization-based systems such as Llama 3 on benchmarks at sizes up to 8B parameters and datasets comprising 4T bytes, and can achieve up to a 50\% reduction in inference computational cost with only marginal impacts on evaluation performance. Additionally, {\textsc BLT} facilitates a novel scaling strategy, permitting concurrent expansion of model capacity and patch size under a constant inference budget—a feature particularly well-suited to the constraints faced in real-world compute environments. By processing raw byte sequences directly, {\textsc BLT} enhances the model’s handling of rare or noisy data instances, thus strengthening its robustness and enriching its ability to capture sub-word-level patterns. Collectively, these findings establish {\textsc BLT} as a compelling alternative to conventional tokenization pipelines, offering a robust, adaptable, and computationally efficient framework for large-scale language modeling.

\end{document}